\section{Conclusion}
\textbf{Capabilities.} The system recognizes six activities on unseen phones with
\FinalAcc\,\% validation accuracy at \FinalMeanMs\,ms per sample on a CPU. It does so by turning
the device-shift problem into a normalization problem: discover devices without labels, then
standardize each with its own statistics. The pipeline is reproducible end to end
(\texttt{poe train}, \texttt{poe experiments}, Docker) and every report number is generated
from code.

\textbf{Limitations.}
(i)~The correction is first-order: it cancels per-feature offsets and gains, not nonlinear
distortions or cross-feature mixing.
(ii)~Per-device statistics assume similar activity mixes across devices; a device dominated by
one activity would be mis-centred (prior shift).
(iii)~Adaptation needs a batch: device groups under 12 records fall back to training statistics
and a logistic-regression fallback (about 88\,\% per file), and folders of 10--20 files remain the
weakest case (86--88\,\%); and outlier \emph{records} contaminate a device's statistics (85\,\% vs.\ 95\,\% for a
global scaler at 20\,\% outliers).
(iv)~Only one labelled training device was available, and validation labels guided model
selection, so the validation score is optimistic for truly new phones.
(v)~Sitting and Standing remain confusable.

\textbf{Future work.} Robust or trimmed per-device statistics against outlier records; second-order alignment (CORAL~\cite{sun2016coral}) or class-balanced
statistics estimated with pseudo-labels to address prior shift; leave-one-device-out
training once Device\_A data is available; a posture-specific second-stage classifier for
Sitting/Standing; and per-sample adaptation for streaming use.
